%# -*- coding: utf-8-unix -*-
\chapter*{ORIGINAL PAPERS}
\addcontentsline{toc}{chapter}{Original Papers}
\markboth{ORIGINAL PAPERS}{ORIGINAL PAPERS}

% hanglist environment is defined globally in trans_number_baker.tex

The following is a list of the principal papers relating to transcendental number\index{transcendental number} theory. It contains most of the recent works on the subject, together with a selection of the more significant of the earlier memoirs. A fuller list of publications prior to 1967 is given in the survey article of Feldman\index{Feldman, N. I.} and Shidlovsky\index{Shidlovsky, A. B.} (see Bibliography).

Abbreviations have been adopted as follows:

\noindent \textbf{C.R.}: Comptes rendus de l'Académie des Sciences (Paris).\\
\noindent \textbf{D.A.N.}: Doklady Akad. Nauk SSSR.\\
\noindent \textbf{I.A.N.}: Izvestia Akad. Nauk SSSR.\\
\noindent \textbf{I.M.}: Indagationes Math.\\
\noindent \textbf{J.M.}: Journal reine angew. Math.\\
\noindent \textbf{M.A.}: Math. Annalen.\\
\noindent \textbf{M.Z.}: Math. Zeitschrift.\\
\noindent \textbf{N.A.W.}: Nederl. Akad. Wetensch. Proc. Ser. A.\\
\noindent \textbf{P.C.P.S.}: Proc. Cambridge Philos. Soc.\\
\noindent \textbf{U.M.}: Uspehi Mat. Nauk.\\
\noindent \textbf{V.M.}: Vestnik Moskov. Univ. Ser. I.\\
I.M. is a collection of the mathematical papers in N.A.W., and reference to the former is given where possible. Mathematical articles in D.A.N. are translated in Soviet Math. Doklady and, where appropriate, reference to the latter is included after an equality sign. A similar convention applies with respect to translations of articles in Trudy Moskov., Mat. Sb., and Mat. Zametki; these are available in Transactions Moscow Math. Soc., Math. USSR Sbornik, and Math. Notes respectively. The abbreviation A.M.S.T. signifies the series American Math. Soc. Translations which contains some articles from other journals. Titles originally in Russian have been given in English throughout.


\begin{hanglist}
    \item ADAMS, W. W. Transcendental numbers\index{transcendental number}\index{numbers} in the p-adic domain, Amer. J. Math 88 (1966), 279-308.
    \item Simultaneous asymptotic diophantine approximations, Mathematika, 14 (1967), 173-80.
    \item Asymptotic diophantine approximations and Hurwitz\index{Hurwitz, A.} numbers, Amer. J. Math. 89 (1967), 1083-108.
    \item Simultaneous asymptotic diophantine approximations to a basis of a real cubic number field, J. Number Th. 1 (1969), 179-94.
    \item ANFERTEVA, E. A. and CHUDAROV, N. G. Minima of the norm function in imaginary quadratic fields, D.A.N. 183 (1968), 255-6; = 9 (1968), 1342-4.
    \item Effective estimates from below of the norms of the ideals of an imaginary quadratic field, Mat. Sb. 82 (1970), 55-66; = 11 (1970), 47-58.
    \item ARMITAGE, J. V. The Thue\index{Thue, A.}-Siegel\index{Siegel, C. L.}-Roth\index{Roth, K. F.} theorem in characteristic p, J. Algebra, 9 (1968), 183-9.
    \item An analogue of a problem of Littlewood, Mathematika, 16 (1969), 101-5; Corrigendum and Addendum, ibid. 17 (1970), 173-8.
    \item Ax\index{Ax, J.}, J. On the units of an algebraic number\index{algebraic number} field, \textit{Illinois J. Math.} 9 (1965), 584-9.
    \item \hspace{1.5em} On Schanuel's conjectures\index{Schanuel's conjecture}, Ann. of Math. 93 (1971), 252-68.
    \item Babaev\index{Babaev, G.}, G. The distribution of integer points on algebraic surfaces, Tadžuk. Gosudarstv. Univ. Dushanbe, 1966.
    \item Baker\index{Baker, R. C.}\index{Baker, A.}, A. Continued fractions\index{continued fractions} of transcendental numbers\index{transcendental number}\index{numbers}, \textit{Mathematika}, 9 (1962), 1-8.
    \item \hspace{1.5em} Rational approximations to certain algebraic numbers\index{algebraic number}, \textit{Proc. London Math. Soc.} 4 (1964), 385-98.
    \item \hspace{1.5em} Rational approximations to \$2 and other algebraic numbers\index{algebraic number}, Quart. J. Math. Oxford, 15 (1964), 375-83.
    \item \hspace{1.5em} On Mahler\index{Mahler, K.}'s classification\index{Mahler's classification} of transcendental numbers\index{transcendental number}\index{numbers}, \textit{Acta Math.} 111 (1964), 97–120.
    \item \hspace{1.5em} On an analogue of Littlewood's Diophantine approximation problem, Michigan Math. J. 11 (1964), 247-50.
    \item \hspace{1.5em} Approximations to the logarithms of certain rational numbers\index{numbers}, Acta Arith. 10 (1964), 315-23.
    \item \hspace{1.5em} Power series representing algebraic functions, J. London Math. Soc. 40 (1965), 108-10.
    \item \hspace{1.5em} On some Diophantine inequalities involving the exponential function\index{exponential function}, Canadian J. Math. 17 (1965), 616-26.
    \item \hspace{1.5em} On a theorem of Sprindžuk\index{Sprindzhuk, V. G.}\index{Sprindžuk, V. G.}, \textit{Proc. Roy. Soc. London}, A 292 (1966), 92-104. A note on the Padé table, \textit{I.M.} 28 (1966), 596-601.
    \item \hspace{1.5em} Linear forms in the logarithms of algebraic numbers\index{logarithms of algebraic numbers}\index{algebraic number} I, II, III, IV, Mathematika, 13 (1966), 204-16; 14 (1967), 102-7, 220-8; 15 (1968), 204-16.
    \item \hspace{1.5em} On Mahler\index{Mahler, K.}'s classification\index{Mahler's classification} of transcendental numbers\index{transcendental number}\index{numbers}: II Simultaneous Diophantine approximation, Acta Arith. 12 (1967), 281-8.
    \item \hspace{1.5em} Simultaneous rational approximations to certain algebraic numbers\index{algebraic number}, P.C.P.S. 63 (1967), 693-702.
    \item \hspace{1.5em} A note on integral integer-valued functions\index{integral integer-valued functions} of several variables, P.C.P.S. 63 (1967), 715-20.
    \item \hspace{1.5em} Contributions to the theory of Diophantine equations\index{Diophantine equations}: I On the representation of integers by binary forms; II The Diophantine equation  $y^2 = x^3 + k$ , Phil. Trans. Roy. Soc. London, A 263 (1968), 173–208.
    \item \hspace{1.5em} The Diophantine equation  $y^2 = ax^3 + bx^2 + cx + d$ , J. London Math. Soc. 43 (1968), 1-9.
    \item \hspace{1.5em} Bounds for the solutions of the hyperelliptic equation\index{hyperelliptic equation}, P.C.P.S. 65 (1969), 439-44.
    \item \hspace{1.5em} A remark on the class number of quadratic fields, Bull. London Math. Soc. 1 (1969), 98-102.
    \item \hspace{1.5em} On the periods of the Weierstrass\index{Weierstrass, K.}  $\wp$ -function, Symposia Math. IV, INDAM Rome, 1968 (Academic Press, London, 1970) pp. 155-74.
    \item \hspace{1.5em} On the quasi-periods of the Weierstrass\index{Weierstrass, K.} \(\zeta\)-function, Göttinger Nachr. (1969) No. 16, 145-57.
    \item \hspace{1.5em} An estimate for the  $\wp$ -function at an algebraic point, Amer. J. Math. 92 (1970), 619-22.
    \item \hspace{1.5em} Effective methods in Diophantine problems I, II, Proc. Symposia Pure Math., vol. 20 (Amer. Math. Soc., 1971), pp. 195-205; ibid. vol. 24, pp. 1-7.
    \item \hspace{1.5em} Effective methods in the theory of numbers\index{numbers}, Proc. Internat. Congress Math. Nice, 1970, vol. 1 (Gauthier-Villars, Paris, 1971), pp. 19-26.
    \item \hspace{1.5em} Imaginary quadratic fields with class number 2, Ann. Math. 94 (1971), 139-52.
    \item BAKER, A. On the class number of imaginary quadratic fields, Bull. Amer. Math. Soc 77 (1971), 678-84.
    \item \hspace{1.5em} A sharpening of the bounds for linear forms in logarithms\index{linear forms in logarithms} I, II, III, Acta Arith. 21 (1972), 117-29; 24 (1973), 33-6; 27 (1975), 247-52.
    \item \hspace{1.5em} A central theorem in transcendence theory, Diophantine approximation and its applications (Academic Press, New York, 1973), pp. 1-23.
    \item BAKER, A., BIRCH, B. J. and WIRSING, E. A. On a problem of Chowla\index{Chowla, S.}\index{Chowla, P.}, J. Number Th. 5 (1973), 224-36.
    \item BAKER, A. and COATES, J. Integer points on curves of genus\index{curves of genus} 1, P.C.P.S. 67 (1970), 595-602.
    \item \textbf{BAKER}, A. and DAVENPORT, H. The equations  $3x^2 2 = y^2$  and  $8x^2 7 = z^2$ , Quart. J. Math. Oxford, Ser. (2), 20 (1969), 129-37.
    \item BAKER, A. and Schinzel\index{Schinzel, A.}, A. On the least integers represented by the genera\index{genera} of binary quadratic forms, \textit{Acta Arith.} 18 (1971), 137-44.
    \item BAKER, A. and SCHMIDT, W. M. Diophantine approximation and Hausdorff dimension\index{Hausdorff dimension}, \textit{Proc. London Math. Soc.} 21 (1970), 1-11.
    \item BAKER, A. and STARK, H. M. On a fundamental inequality in number theory, Ann. Math. 94 (1971), 190-9.
    \item BALKEMA, A. A. and TIJDEMAN, R. Some estimates in the theory of exponential sums, Acta Math. Hungar. 24 (1973), 115-133.
    \item BARON, G. and BRAUNE, E. Zur Transzendenz von Luckenriehen mit ganzalgebraischen Koeffizienten und algebraischem Argument, Compositio Math. 22 (1970), 1-6, and a similar title, Arch. Math. (Basel) 22 (1971), 379-84.
    \item BELOGRIVOV, I. I. Transcendence and algebraic independence of values of certain E-functions\index{E-functions}, V.M. 22 (1967), 55-62, and similar titles, D.A.N. 174 (1967), 267-70; = 8 (1967), 610-3; Mat. Sb. 82 (1970), 387-408; Sibirsk Mat. Z. 12 (1971), 961-82; 14 (1973), 16-35.
    \item BOEHLE, K. Über die Transzendenz von Potenzen mit algebraischen Exponenten (Verallgemeinerung eines Satzes von A. Gelfond\index{Gelfond, A. O.}), M.A. 108 (1933), 56-74. Über die Approximation von Potenzen mit algebraischen Exponenten durch algebraische Zahlen, M.A. 110 (1935), 662-78.
    \item Bombieri\index{Bombieri, E.}, E. Sull'approssimazione di numeri algebrici mediante numeri algebrici, Boll. Un. Mat. Ital. 13 (1958), 351-4.
    \item \hspace{1.5em} Algebraic values of meromorphic maps, \textit{Invent. Math.} 10 (1970), 267-87; addendum, \textit{ibid.} 11 (1970), 163-6.
    \item Bombieri\index{Bombieri, E.}, E. and Lang\index{Lang, S.}, S. Analytic subgroups of group varieties, \textit{Invent.}Math. 11 (1970), 1-14.
    \item Borel\index{Borel, E.}, E. Sur la nature arithmétique du nombre e, C.R. 128 (1899), 596 9.
    \item BOYD, D. W. Transcendental numbers\index{transcendental number}\index{numbers} with badly distributed powers, \textit{Proc. Amer. Math. Soc.} 23 (1969), 424-7.
    \item Brauer\index{Brauer, R.}\index{Brauer, A.}, A. Über diophantische Gleichungen mit endlich vielen Lösungen, J.M. 160 (1929), 70-99; 161 (1929), 1-13.
    \item Brownawell\index{Brownawell, W.}, W. Some transcendence results for the exponential function\index{exponential function}, Norske Vid. Selsk. Skr. (1972), no. 11.
    \item \hspace{1.5em} The algebraic independence of certain values of the exponential function\index{exponential function}, \textit{ibid.} (1972), no. 23, and a similar title, J. Number Th. 6 (1974), 22-31.
    \item \hspace{1.5em} Sequences of Diophantine approximations, J. Number Th. 6 (1974), 11-21.
    \item BRUMER, A. On the units of algebraic number\index{algebraic number} fields, Mathematika, 14 (1967), 121-4.
    \item Bundschuh\index{Bundschuh, P.}, P. Ein Satz über ganze Funktionen und Irrationalitätaussagen, \textit{Invent. Math.} 9 (1969), 175–84.
    \item Ein Approximationsmass für transzendente Lösungen gewisser transzendenter Gleichungen, J.M. 251 (1971), 32-53.
    \item Irrationalitätsmasse für  $e^a$ ,  $a \neq 0$ , rational oder Liouville\index{Liouville, J.} Zahl, M.A. 192 (1971), 229-42.
    \item Bundschur, P. and Hock\index{Hock, A.}, A. Bestimmung aller imaginär-quadratischen Zahlkörper der Klassenzahl Eins mit Hilfe eines Satzes von Baker\index{Baker, R. C.}\index{Baker, A.}, M.Z. 111 (1969), 191-204.
    \item CANTOR, G. Über eine Eigenschaft des Inbegriffs aller reellen algebraischem Zahlen, J.M. 77 (1874), 258-62; = Ges. Abh. 116-8.
    \item Chowla\index{Chowla, S.}\index{Chowla, P.}, P. Remarks on a previous paper, J. Number Th. 1 (1969), 522-4.
    \item CHUDAKOV, N. G. The upper bound for the discriminant of the tenth imaginary quadratic field with class number 1, Studies in number theory, no. 3. Saratov (1969), 73-7.
    \item Chudnovsky\index{Chudnovsky, G. V.}, G. V. Algebraic independence of several values of the exponential function\index{exponential function}, \textit{Mat. Zametki}, 15 (1974), 661-72.
    \item CIJSOUW, P. L. Transcendence measures (dissertation, Amsterdam, 1972).
    \item CIJSOUW, P. L. and TIJDEMAN, R. Distinct prime factors of consecutive integers, Diophantine approximation and its applications (Academic Press, London, 1973), pp. 59-76.
    \item \hspace{1.5em} Transcendence of certain power series of algebraic numbers\index{algebraic number}, \textit{Acta Arith.} 23 (1973), 301-5.
    \item COATES, J. On the algebraic approximation of functions I, II, III, IV, I.M. 28 (1966), 421-61; 29 (1967), 205-12.
    \item \hspace{1.5em} An effective p-adic analogue of a theorem of Thue\index{Thue, A.} I; II: The greatest prime factor of a binary form; III: The Diophantine equation  $y^2 = x^3 + k$ , Acta Arith. 15 (1969), 279-305; 16 (1970), 399-412, 425-35.
    \item \hspace{1.5em} Construction of rational functions on a curve, P.C.P.S. 68 (1970), 105-23.
    \item \hspace{1.5em} An application of the Thue\index{Thue, A.}-Siegel\index{Siegel, C. L.}-Roth\index{Roth, K. F.} theorem to elliptic functions\index{elliptic functions}, P.C.P.S. 69 (1971), 157-61.
    \item \hspace{1.5em} Linear forms in the periods of the exponential and elliptic functions\index{elliptic functions}, \textit{Invent.}Math. 12 (1971), 290-9.
    \item \hspace{1.5em} The transcendence of linear forms in  $\omega_1$ ,  $\omega_2$ ,  $\eta_1$ ,  $\eta_2$ ,  $2\pi i$ , Amer. J. Math. 93 (1971), 385-97.
    \item \hspace{1.5em} Linear relations between  $2\pi i$  and the periods of two elliptic curves\index{elliptic curve}, Diophantine approximation and its applications (Academic Press, London 1973), pp. 77-99.
    \item Cugiani\index{Cugiani, M.}, M. Sulla approssimabilità dei numeri algebrici mediante numeri razionali, Ann. Mat. Pura Appl. 48 (1959), 135-45, and a similar title, Boll. Un. Mat. Ital. 14 (1959), 151-62.
    \item DAVENPORT, H. A note on binary cubic forms, Mathematika, 8 (1961), 58-62. A note on Thue\index{Thue, A.}'s theorem, ibid. 15 (1968), 76-87.
    \item DAVENPORT, H. and ROTH, K. F. Rational approximations to algebraic numbers\index{rational approximations to algebraic numbers}\index{algebraic number}, \textit{Mathematika}, 2 (1955), 160-7.
    \item DAVENPORT, H. and SCHMIDT, W. M. Approximation of real numbers by algebraic integers\index{algebraic integer}, \textit{Acta Arith.} 15 (1969), 393-416.
    \item DEMJANENKO, V. A. The representation of numbers\index{numbers} by a binary cubic irreducible form, Mat. Zametki, 7 (1970), 87-96.
    \item Dyson\index{Dyson, F.}, F. J. The approximation to algebraic numbers\index{algebraic number} by rationals, \textit{Acta Math.} 79 (1947), 225-40.
    \item ELLISON, W. J., PESEK, J., STALL, D. S. and LUNNON, W. F. A postscript to a paper of A. Baker\index{Baker, R. C.}\index{Baker, A.}, Bull. London Math. Soc. 3 (1971), 75-8.
    \item ELLISON, W. J., ELLISON, F., PESEK, J., STAHL, C. E. and STALL, D. S. The Diophantine equation y\ + k = x\\textit{, J. Number Th.} 4 (1972), 107-17.
    \item ERDOS, P. Representations of real numbers as sums and products of Liouville\index{Liouville, J.} numbers\index{Liouville number}, \textit{Bull. Amer. Math. Soc.} 68 (1962), 475-8.
    \item ERDOS, P., and MAHLER, K. Some arithmetical properties of the convergents of a continued fraction, \textit{J. London Math. Soc.} 14 (1939), 12-18.
    \item FADDEEV, D. K. On a paper by A. Baker\index{Baker, R. C.}\index{Baker, A.}, \textit{Zap. Nau6n. Sem. Leningrad Otdel. Mat. Inst. Steklov,} 1 (1966), 128-39.
    \item FELDMAN, N. I. The approximation of some transcendental numbers\index{transcendental number}\index{numbers} I", If, \textit{I.A.N.} 15 (1951), 53-74; = \textit{A.M.S.T.} 59 (1969), 224-70.
    \item \hspace{1.5em} On the simultaneous approximation of the periods of elliptic functions\index{elliptic functions} by algebraic numbers\index{algebraic number}, \textit{I.A.N.} 22 (1958), 563-76; = \textit{A.M.S.T.} 59 (1966), 271-84.
    \item \hspace{1.5em} On the measure of transcendence\index{measure of transcendence} of \(\pi\) 24(1960), 357-68; = \textit{A.M.S.T.} 58(1966), 110-24.
    \item \hspace{1.5em} On the approximation by algebraic numbers\index{algebraic number} of the logarithms of algebraic numbers\index{logarithms of algebraic numbers}\index{algebraic number}, \textit{I.A.N.} 24 (1960), 475-92; = \textit{A.M.S.T.} 58 (1966), 125-42.
    \item \hspace{1.5em} On transcendental numbers\index{transcendental number}\index{numbers} having approximations of a given type, \textit{U.M.} 17 (1962), 145-51.
    \item \hspace{1.5em} On the problem of the measure of transcendence\index{measure of transcendence} of e, \textit{U.M.} 18 (1963), 207-13.
    \item \hspace{1.5em} Arithmetic properties of the solutions of some transcendental equations, \textit{V.M.} 1 (1964), 13-20.
    \item \hspace{1.5em} An estimate of the absolute value of a linear form in the logarithms of certain algebraic numbers\index{algebraic number}, \textit{Mat. Zametki, 2} (1967), 245-56, and a similar title, \textit{V.M.} 2 (1967), 63-72.
    \item \hspace{1.5em} Estimation of a linear form in the logarithms of algebraic numbers\index{logarithms of algebraic numbers}\index{algebraic number}, \textit{Mat. Sb.} 76 (1968), 304-19; = 5 (1968), 291-307, and similar titles, \textit{Mat. Sb.} 77 (1968), 423-36; = 6 (1968), 393-406; \textit{U.M.} 23 (1968), 185-6; \textit{D.A.N.} \textbf{182} (1968), 1278-9; = 9 (1968), 1284-5.
    \item \hspace{1.5em} An elliptic analogue of an inequality of A. O. Gelfond\index{Gelfond, A. O.}, \textit{Trudy Moskov.} 18 (1968), 65-76; = 18 (1968), 71-83.
    \item \hspace{1.5em} Refinement of two effective inequalities of A. Baker\index{Baker, R. C.}\index{Baker, A.}, \textit{Mat. Zametki,} 6 (1969), 767-9, and related work, 5 (1969), 681-9.
    \item \hspace{1.5em} Effective bounds for the size of the solutions of certain Diophantine equations\index{Diophantine equations}, \textit{Mat. Zametki,} 8 (1970), 361-71; = 8 (1970), 674-9, and related work, \textit{ibid. 7} (1970), 569-80; = 7 (1970), 343-9; \textit{V.M.} 26 (1971), 52-8.
    \item \hspace{1.5em} An effective refinement of the exponent in Liouville\index{Liouville, J.}'s theorem\index{Liouville's theorem}, \textit{I.A.N.} 35 (1971), 973-90; = \textit{Math. USSR Izv.} 5 (1971), 985-1002; and related work \textit{D.A.N.} 207 (1972), 41-3.
    \item \hspace{1.5em} On a certain linear form, \textit{Acta Arith.} 21 (1972), 347-55.
    \item FELDMAN, N. I. and CHUDAKOV, N.G. On a theorem of Stark\index{Stark, H. M.}, \textit{Mat. Zametki,} 3 (1972), 329-40.
    \item FRAENKEL, A. S. On a theorem of D. Ridout\index{Ridout, D.} in the theory of Diophantine approximations, \textit{Trans. Amer. Math. Soc.} \textbf{105} (1962), 84-101.
    \item \hspace{1.5em} Transcendental numbers\index{transcendental number} and a conjecture of Erdos\index{Erdos, P.} and Mahler\index{Mahler, K.}, \textit{J. London Math. Soc.} 39 (1964), 405-16.
    \item \hspace{1.5em} Distance to the nearest integer and algebraic independence of certain real numbers\index{numbers}, \textit{Proc. Amer. Math. Soc.} 16 (1965), 154-60.
    \item FRAENKEL, A. S.andBoROSH, I. Fractional dimension of a set of transcendental numbers\index{transcendental number}\index{numbers}, \textit{Proc. London Math. Soc.} 15 (1965), 458-70.
    \item GALOCKIN, A. I. Bounds for the measure of transcendence\index{measure of transcendence} of the values of E-functions\index{E-functions}, \textit{Mat. Zametki,} 3 (1968), 377-86, and similar titles, \textit{ibid.} 8 (1970), 19-28; = 8 (1970), 478-84; \textit{V.M.} 25 (1970), 58-63.
    \item GELFOND, A.O. Sur les propriétés arithmetiques des fonctions entières, Tôhoku Math. J. 30 (1929), 280-5.
    \item \hspace{1.5em} Sur les nombres transcendants, C.R. 189 (1929), 1224-6.
    \item \hspace{1.5em} Sur le septième problème de Hilbert\index{Hilbert, D.}, I.A.N. 7 (1934), 623-30; D.A.N. 2 (1934), 1-6.
    \item \hspace{1.5em} On the approximation of transcendental numbers\index{transcendental number} by algebraic numbers\index{algebraic number}, D.A.N. 2 (1935), 177-82.
    \item \hspace{1.5em} On the approximation by algebraic numbers\index{algebraic number} of the ratio of the logarithms of two algebraic numbers\index{algebraic number}, I.A.N. 5-6 (1939), 509-18.
    \item \hspace{1.5em} Sur la divisibilité de la différence des puissances de deux nombres entières par une puissance d'un idéal premier, Mat. Sb. 7 (1940), 7-26.
    \item \hspace{1.5em} The approximation of algebraic numbers\index{algebraic number} by algebraic numbers\index{algebraic number} and the theory of transcendental numbers\index{transcendental number}\index{numbers}, U.M. 4 (1949), 19-49; = A.M.S.T. 2 (1962), 81-124.
    \item \hspace{1.5em} On the algebraic independence of algebraic powers of algebraic numbers\index{algebraic number}, D.A.N. 64 (1949), 277-80, and similar titles, D.A.N. 67 (1949), 13-14; U.M. 5 (1949), 14-48; = A.M.S.T. 2 (1962), 125-69.
    \item GELFOND, A. O. and FELDMAN, N. I. On lower bounds for linear forms in three logarithms of algebraic numbers\index{logarithms of algebraic numbers}\index{algebraic number}, V.M. 5 (1949), 13-16.
    \item \hspace{1.5em} On the measure of mutual transcendence of certain numbers\index{numbers}, I.A.N. 14 (1950), 493-500.
    \item GELFOND, A. O. and LINNIK, Yu. V. On Thue\index{Thue, A.}'s method and the problem of effectiveness in quadratic fields, D.A.N. 61 (1948), 773-6.
    \item GOLDSTEIN, L. J. A generalization of Stark\index{Stark, H. M.}'s theorem, J. Number Th. 3 (1971), 323-46.
    \item \hspace{1.5em} Imaginary quadratic fields of class number 2, ibid. 4 (1972), 286-301.
    \item GORDAN, P. Transzendenz von e und  $\pi$ , M.A. 43 (1893), 222-5.
    \item GÜNTHER\index{Günther, A.}, A. Über transzendente φ-adische Zahlen I, II, J.M. 192 (1953), 155-66; 193 (1954), 1-10.
    \item GÜTING, R. Approximation of algebraic numbers\index{algebraic number} by algebraic numbers\index{algebraic number}, \textit{Michigan Math. J.} 8 (1961), 149-59.
    \item \hspace{1.5em} On Mahler\index{Mahler, K.}'s function  $\theta_1$ , \textit{ibid.} 10 (1963), 161-79.
    \item \hspace{1.5em} Uber die Zusammenhang zwischen rationalen Approximationen und Kettenbrüchentwicklungen, M.Z. 90 (1965), 382-7.
    \item Polynomials with multiple zeros, Mathematika, 14 (1967), 181-96.
    \item HANECKE, W. Über ein System simultaner diophantischer Approximationen, M.Z. 110 (1969), 378-84.
    \item HERMITE, CH. Sur la fonction exponentielle, C.R. 77 (1873), 18-24, 74-9, 226-33, 285-93; = Oeuvres III, 150-81.
    \item HILBERT, D. Über die Transzendenz der Zahlen e und  $\pi$ , Göttinger Nachr. (1893), 113-16; M.A. 43 (1893), 216-20; = Ges. Abh. I, 1-4.
    \item HILLIKER, D. L. On analytic functions which have algebraic values at a convergent sequence of points, Trans. Amer. Math. Soc. 126 (1967), 534-50.
    \item HURWITZ, A. Ueber arithmetische Eigenschaften gewisser transcendenter Functionen I, II, M.A. 22 (1883), 211-29; 32 (1888), 583-8.
    \item \hspace{1.5em} Beweis der Transzendenz der Zahl e, Göttinger Nachr. (1893), 153-5.
    \item HYYRÖ, S. Über Approximation algebraischer Zahlen durch rationale, Ann. Univ. Turku. ser. AI, 84 (1965), 1-12, and similar titles, Ann. Acad. Sci. Fenn. ser. AI, nos. 376 (1965), 394 (1967).
    \item IÇEN, O. S. Eine Verallgemeinerung und Übertragung der Schneiderschen Algebraizitätskriterien ins p-adische mit Anwendung auf einen Transzendenzbeweis im p-adischen, J.M. 198 (1957), 28-55; Berichtigung, ibid. 220.
    \item KAPPE, LUISE-CHARLOTTE\index{Kappe, Luise-Charlotte}. Zur Approximation von ea, Ann. Univ. Sci. Budapest Eötvös Sect. Math. 9 (1966), 3-14.
    \item KASCH, F. Zur Annäherung algebraischer Zahlen durch arithmetisch charakterisierte rationale Zahlen, \textit{Math. Nachr.} 10 (1953), 85–98.
    \item \hspace{1.5em} Über eine metrische Eigenschaft der S-Zahlen, M.Z. 70 (1958), 263-70, and a similar title, J.M. 203 (1960), 157-9.
    \item KASCH, F. and VOLKMANN, B. Zur Mahlerschen Vermutung über S-Zahlen, M.A. 136 (1958), 442-53.
    \item KAUFMAN, R. M. Bounds for linear forms in the logarithms of algebraic numbers\index{logarithms of algebraic numbers}\index{algebraic number} with p-adic metric, V.M. 26 (1971), 3-10.
    \item KEATES, M. On the greatest prime factor of a polynomial, \textit{Proc. Edinburgh Math. Soc.} 16 (1969), 301-3.
    \item Koksma\index{Koksma, J. F.}, J. Über die Mahlersche Klasseneinteilung der transzendenten Zahlen und die Approximation komplexer Zahlen durch algebraische Zahlen, Monatsh. Math. Phys. 48 (1939), 176–89.
    \item Korsma, J. F. and Popken\index{Popken, J.}, J. Zur Transzendenz von  $e^{\pi}$\index{transcendence of  $e^{\pi}$} , J.M. 168 (1932), 211-30.
    \item Kuzmin\index{Kuzmin, R. O.}, R. O. On a new class of transcendental numbers\index{transcendental number}\index{numbers}, I.A.N. 3 (1930), 583-97.
    \item LAMBERT, J. H. Mémoire sur quelques propriétés remarquables des quantités transcendantes circulaires et logarithmiques, \textit{Histoire Acad. roy. sci. et belles lettr.} Berlin, Année 1761 (1768), 265–322.
    \item LANG, S. On a theorem of Mahler\index{Mahler, K.}, Mathematika, 7 (1960), 139-40.
    \item \hspace{1.5em} A transcendence measure for E-functions\index{E-functions}, ibid. 9 (1962), 157-61.
    \item \hspace{1.5em} Transcendental points on group varieties, Topology, 1 (1962), 313-8.
    \item \hspace{1.5em} Asymptotic approximations to quadratic irrationalities I, II, Amer. J. Math. 87 (1965), 481-95.
    \item \hspace{1.5em} Report on Diophantine approximations, Bull. Soc. Math. France, 93 (1965), 177-92.
    \item \hspace{1.5em} Algebraic values of meromorphic maps I, II, \textit{Topology}, 3 (1965), 183-91; 5 (1966), 363-70.
    \item \hspace{1.5em} Transcendental numbers\index{transcendental number}\index{numbers} and Diophantine approximations, Bull. Amer. Math. Soc. 77 (1971), 635-77.
    \item LEVEQUE, W. J. Note on the transcendence of certain series, \textit{Proc. Amer. Math. Soc.} 2 (1951), 401-3.
    \item \hspace{1.5em} Note on S-numbers\index{S-number}\index{S*-number}\index{numbers}, Proc. Amer. Math. Soc. 4 (1953), 189-90.
    \item \hspace{1.5em} On Mahler\index{Mahler, K.}'s U-numbers\index{U-number}\index{U*-number}, J. London Math. Soc. 28 (1953), 220-9.
    \item LEVIN, B. V., FELDMAN, N. I. and SHIDLOVSKY, A. B. Alexander O. Gelfond\index{Gelfond, A. O.}, Acta Arith. 17 (1971), 315-36.
    \item LINDEMANN, F. Über die Zahl  $\pi$ , M.A. 20 (1882), 213-25.
    \item \hspace{1.5em} Über die Ludolph'sche Zahl, S.-B. Preuss Akad. Wiss. (1882), 679-82.
    \item LINNIK, YU. V. and CHUDAKOV, N. G. On a class of completely multiplicative functions, D.A.N. 74 (1950), 193-6.
    \item LIOUVILLE, J. Sur l'irrationalité du nombre e, J. Math. pures appl. 5 (1840), 192-3.
    \item \hspace{1.5em} Sur des classes très-étendues de quantités dont la valeur n'est ni algébrique, ni même reductible à des irrationnelles algébriques, C.R. 18 (1844), 883-5, 910-11, J. Math. pures appl. 16 (1851), 133-42.
    \item MAHLER, K. Über Beziehungen zwischen der Zahl e und den Liouvilleschen Zahlen, M.Z. 31 (1930), 729-32.
    \item \hspace{1.5em} Ein Beweis des Thue\index{Thue, A.}-Siegelschen Satzes über die Approximation algebraischen Zahlen für binomische Gleichungen, M.A. 105 (1931), 267-76.
    \item Zur Approximation der Exponentialfunktion und des Logarithmus I, II, J.M. 166 (1932), 118-50.
    \item Ein Beweis der Transzendenz der p-adischen Exponentialfunktion, J.M. 169 (1932), 61-6.
    \item Über das Mass der Menge aller S-Zahlen, M.A. 106 (1932), 131-9.
    \item Zur Approximation algebraischer Zahlen: I Über den grössten Primteiler binärer Formen, II Über die Anzahl der Darstellungen grosser Zahlen durch binäre Formen, III Über die mittlere Anzahl der Darstellungen grosser Zahlen durch binäre Formen, M.A. 107 (1933), 691-730; 108 (1933), 37-55; Acta Math. 62 (1933), 91-166.
    \item Uber transzendente p-adische Zahlen, Compositio Math. 2 (1935), 259-75, and related work, Mathematica (Leiden), 3 (1935), 177-85.
    \item Ein Analogon zu einen Schneiderschen Satz, N.A.W. 39 (1936), 633-40, 729-37.
    \item Arithmetische Eigenschaften einer Klasse von Dezimalbrüchen, N.A.W. 40 (1937), 421-8, and related work, Math. B. Zutphen, 6 (1937), 22-36.
    \item On the fractional parts of the powers of a rational number I, II, Acta Arith. 3 (1938), 89-93; Mathematika, 4 (1957), 122-4.
    \item On the solution of algebraic differential equations, N.A.W. 42 (1939), 61-3. On the theorem of Liouville\index{Liouville, J.} in fields of positive characteristic, Canadian J. Math. 1 (1949), 397-400.
    \item On Dyson\index{Dyson, F.}'s improvement of the Thue\index{Thue, A.}-Siegel\index{Siegel, C. L.} theorem, I.M. 11 (1949), 449–58.
    \item On the continued fractions\index{continued fractions} of quadratic and cubic irrationals, Ann. Math. pura appl. 30 (1949), 147-72.
    \item On the generating functions of the integers with a missing digit, J. Indian Math. Soc. A 15 (1951), 33-40.
    \item On the greatest prime factor of  $ax^m + by^n$ , N.A.W. 1 (1953), 113-22.
    \item On the approximation of logarithms of algebraic numbers\index{logarithms of algebraic numbers}\index{algebraic number}, \textit{Phil. Trans. Roy. Soc. London}, A 245 (1953), 371-98.
    \item On the approximation of  $\pi$ , I.M. 15 (1953), 30-42.
    \item On a theorem of E. Bombieri\index{Bombieri, E.}, I.M. 22 (1960), 245-53; Correction, ibid. 23 (1961), 141.
    \item An application of Jensen's formula to polynomials, \textit{Mathematika}, 7 (1960), 98-100.
    \item On some inequalities for polynomials in several variables, J. London Math. Soc. 37 (1962), 341-4.
    \item On the approximation of algebraic numbers\index{algebraic number} by algebraic integers\index{algebraic integer}, J. Australian Math. Soc. 1 (1963), 408-34.
    \item Transcendental numbers\index{transcendental number}\index{numbers}, ibid. 4 (1964), 393-6.
    \item An inequality for a pair of polynomials that are relatively prime, \textit{ibid.} 4 (1964), 418-20.
    \item Arithmetic properties of lacunary power series with integral coefficients, ibid. 5 (1965), 56-64.
    \item Applications of some formulae by Hermite\index{Hermite, Ch.} to the approximation of exponentials and logarithms, M.A. 168 (1967), 200-27.
    \item On a class of entire functions, Acta Math. Hungar. 18 (1967), 83-96.
    \item A remark on recursive sequences, J. Math. Sci. 1 (1966), 12-17.
    \item Applications of a theorem of A. B. Shidlovsky\index{Shidlovsky, A. B.}, Proc. Roy. Soc. London, A 305 (1968), 149-73, and a similar title, Mat. Zametki, 2 (1967), 25-32. Perfect systems, Compositio Math. 19 (1968), 95-166.
    \item Remarks on a paper by W. Schwarz\index{Schwarz, W.}, J. Number Th. 1 (1969), 512-21.
    \item MAHLER, K. On the order function of a transcendental number\index{transcendental number}, Acta Arith. 18 (1971), 63-76.
    \item MAHLER, K. and POPKEN, J. Ein neues Prinzip für Transzendenzbeweise N.A.W. 38 (1935), 864-71.
    \item MAHLER, K. and SZEKERES, G. On the approximation of real numbers\index{numbers} by roots of integers, \textit{Acta Arith.} 12 (1967), 315-20.
    \item MAIER, W. Potenzreihen irrationalen Grenzwertes, J.M. 156 (1927), 93-148.
    \item MAILLET, E. Sur les racines des équations transcendantes à coefficients rationnels, \textit{J. Math. pures appl.} 7 (1901), 419-40; \textit{C.R.} 132 (1901), 908-10; 133 (1901), 989-90.
    \item \hspace{1.5em} Sur les propriétés arithmétiques des fonctions entières et quasi entières, C.R. 134 (1902), 1131-2; Bull. Soc. Math. France, 30 (1902), 134-53.
    \item \hspace{1.5em} Sur les fonctions monodromes et les nombres transcendants, J. Math. pures appl. 10 (1904), 275-362; C.R. 138 (1904), 262-5.
    \item \hspace{1.5em} Sur les nombres e et  $\pi$ , et les equations transcendants, C.R. 133 (1901), 1191-2; Acta Math. 29 (1905), 295-331.
    \item \hspace{1.5em} Sur la classification des irrationnelles, C.R. 143 (1906), 26-8.
    \item \hspace{1.5em} Sur les nombres transcendants dont le développement en fraction continue est quasi-périodique, et sur les nombres de Liouville\index{Liouville, J.}, Bull. Soc. Math. France, 34 (1906), 213-27, and related work ibid. 35 (1907), 27-47; 50 (1922), 74-99; J. Math. pures appl. 3 (1907), 299-336; C.R. 170 (1920), 983-6.
    \item MASSER, D. W. On the periods of the exponential and elliptic functions\index{elliptic functions}, P.C.P.S. 73 (1973), 339-50.
    \item \hspace{1.5em} Elliptic functions\index{elliptic functions} and transcendence, Ph.D. dissertation (Cambridge, 1974).
    \item MENDÈS FRANCE, M. Nombres transcendants et ensembles normaux, Acta Arith. 15 (1969), 189-92.
    \item MEYER, Y. Nombres algébriques, nombres transcendants et equirépartition modulo 1, Acta Arith. 16 (1970), 347-50.
    \item Morduchai-Boltovskoj, D. Sur le logarithme d'un nombre algébrique, C.R. 176 (1923), 724-7.
    \item \hspace{1.5em} On some properties of transcendental numbers\index{transcendental number}\index{numbers} of the second class, \textit{Mat. Sb.} \textbf{34} (1927), 55–100.
    \item \hspace{1.5em} On transcendental numbers\index{transcendental number}\index{numbers} with successive approximations defined by algebraic equations, \textit{ibid.} 41 (1934), 221-32.
    \item \hspace{1.5em} Über einige Eigenschaften der transzendenten Zahlen, Töhoku Math. J. 40 (1935), 99-127.
    \item \hspace{1.5em} On conditions for the definition of a number by transcendental equations of some general type D.A.N. 52 (1946), 487-90.
    \item \hspace{1.5em} On hypertranscendental functions and hypertranscendental numbers\index{numbers}, D.A.N. 64 (1949), 21-4.
    \item MOBITA, Y. On transcendency of special values of arithmetic automorphic functions, J. Math. Soc. Japan 24 (1972), 268-74.
    \item NESTERENKO, Ju. V. The algebraic independence of the values of \textit{E}-functions which satisfy linear homogeneous differential equations, \textit{Mat. Zametki} 5 (1969), 587-98.
    \item VON NEUMANN, J. Ein System algebraisch unabhängiger Zahlen, M.A. 99 (1928), 134-41.
    \item Nurmagomedov\index{Nurmagomedov, M. S.}, M. S. The arithmetic properties of a certain class of analytic functions, \textit{Mat. Sb.} 85 (1971), 339-65, and similar titles, \textit{V.M.} 26, no. 6 (1971), 79-86; 28, no. 1 (1973), 19-25.
    \item OLEINIKOV, V. A. On some properties of algebraically dependent quantities, \textit{V.M.} 5 (1962), 11-17.
    \item \hspace{1.5em} On the transcendence and algebraic independence of the values of some E-functions\index{E-functions}, \textit{V.M.} 6 (1962), 34-8.
    \item \hspace{1.5em} On the algebraic independence of the values of E-functions\index{E-functions} that are solutions of third-order linear inhomogeneous differential equations, \textit{D.A.N.} 169 (1966), 32-4; = 7 (1966), 869-71, and related work, \textit{ibid.} 166 (1966), 540-3; \textit{Mat. Sb.} 78 (1969), 301-6.
    \item OSOOOD, C. F. Some theorems on diophantine approximation, \textit{Trans. Amer. Math. Soc.} 123 (1966), 64-87.
    \item \hspace{1.5em} A method in diophantine approximation I-V, \textit{Acta Arith.} 12 (1967), 111-29; 13 (1968), 383-93; 16 (1970), 5-40; 22 (1973), 353-69.
    \item \hspace{1.5em} The simultaneous diophantine approximation of certain \textit{kth} roots, \textit{P.C.P.S.} 67 (1970), 75-86.
    \item \hspace{1.5em} On the diophantine approximation of values of functions satisfying certain $q$-difference equations, \textit{J. Number Th.} 3 (1971), 159-77.
    \item \hspace{1.5em} On the simultaneous diophantine approximation of values of certain algebraic functions, \textit{Acta Arith.} 19 (1971), 343-86.
    \item \hspace{1.5em} A Diophantine property of \textit{Ja(z), Diophantine approximation and its applications} (Academic Press, London, 1973), pp. 201-9.
    \item \hspace{1.5em} An effective lower bound on the ' Diophantine approximation' of algebraic functions by rational functions, \textit{Mathematika,} 20 (1973), 4-15.
    \item PABRY\index{Parry, C. J.}, C. J. The jj-adic generalisation of the Thue\index{Thue, A.}-Siegel\index{Siegel, C. L.} theorem, \textit{J. London Math. Soc.} 15 (1940), 293-305, and a similar title, \textit{Acta Math.} 83 (1950), 1-100.
    \item PECK, L. G. Simultaneous rational approximations to algebraic numbers\index{rational approximations to algebraic numbers}\index{algebraic number}, \textit{Bull. Amer. Math. Soc.} 67 (1961), 197-201.
    \item PERSHIKOVA, T. V. On the transcendence of values of certain J5-functions, \textit{V.M.} 21 (1966), 55-61.
    \item P6LYA\index{Polya, G.}, G. tTber ganzwertige ganze Funktionen, \textit{Rend. Circ. Mat. Palermo,} 40 (1915), 1-16.
    \item VAN DER POORTEN, A. J. Transcendental entire functions mapping every algebraic number\index{algebraic number} field into itself, \textit{J. Australian Math. Soc.} 8 (1968), 192-3.
    \item \hspace{1.5em} On the arithmetic nature of definite integrals of rational functions, \textit{Proc. Amer. Math. Soc.} 29 (1971), 451-6.
    \item POPKEN, J. Sur la nature arithmétique du nombre e, \textit{C.R.} \textbf{186} (1928), 1505-7. Zur Transzendenz von e, \textit{M.Z.} 29 (1929), 525-41.
    \item \hspace{1.5em} Zur Transzendenz von \textit{n, M.Z.} 29 (1929), 542-8.
    \item \hspace{1.5em} Eine arithmetische Eigenschaft gewisser ganzer Funktionen I, II, \textit{N.A.W.} 40 (1937), 142-50, 263-70.
    \item \hspace{1.5em} On Lambert\index{Lambert, J. H.}'s proof for the irrationality of \textit{n, N.A.W.} 43 (1940), 712-14; 52 (1949), 504.
    \item \hspace{1.5em} An arithmetical theorem concerning linear differential equations of infinite order, \textit{I.M.} 12 (1950), 1645-6.
    \item \hspace{1.5em} Un theoreme sur lea nombres transcendants, \textit{Bull. Soc. Math. Belg.} 7 (1955), 124-30.
    \item \hspace{1.5em} Arithmetical properties of the Taylor coefficients of algebraic functions, \textit{I.M.} 21 (1959), 202-10.
    \item \hspace{1.5em} Irrational power series, \textit{I.M.} 25 (1963), 691-4.
    \item \hspace{1.5em} Note on a generalization of a problem of Hilbert\index{Hilbert, D.}, \textit{I.M.} 27 (1965), 178-81.
    \item \hspace{1.5em} Algebraic independence of certain zSta functions, \textit{I.M.} 28 (1966), 1-5.
    \item \hspace{1.5em} A measure for the differential transcendence of the zeta-function of Riemann,
    \item Number theory and analysis (Papers in honour of Edmund Landau\index{Landau, E.}), (Plenum, New York, 1969), pp. 245-55.
    \item POPEEN, J. A contribution to the Thue\index{Thue, A.}-Siegel\index{Siegel, C. L.}-Roth\index{Roth, K. F.} problem, Number theory (North-Holland, Amsterdam, 1970), pp. 181-90.
    \item RAMACHANDRA, K. Approximation of algebraic numbers\index{algebraic number}, Göttinger Nachr. (1966), 45-52.
    \item \hspace{1.5em} Contributions to the theory of transcendental numbers\index{transcendental number}\index{numbers} I, II, Acta Arith. 14 (1968), 65-88.
    \item \hspace{1.5em} A note on Baker\index{Baker, R. C.}\index{Baker, A.}'s method, J. Australian Math. Soc. 10 (1969), 197-203.
    \item \hspace{1.5em} A lattice point problem for norm forms in several variables, J. Number Th. 1 (1969), 534-55.
    \item \hspace{1.5em} A note on numbers\index{numbers} with a large prime factor I, II, III, J. London Math. Soc. 1 (1969), 303-6; J. Indian Math. Soc. 34 (1970), 39-48; Acta Arith. 19 (1971), 49-62.
    \item \hspace{1.5em} Application of Baker\index{Baker, R. C.}\index{Baker, A.}'s theory to two problems considered by Erdös and Selfridge, J. Indian Math. Soc. 37 (1973) 25-34.
    \item RAMACHANDRA, K. and Shorey\index{Shorey, T. N.}, T. N. On gaps between numbers\index{numbers} with a large prime factor, \textit{Acta Arith.} 24 (1973), 99-111.
    \item RAUEY, G. Algébricité des fonctions méromorphes prenant certaines valeurs algébriques, Bull. Soc. Math. France, 96 (1968), 197-208.
    \item RIDOUT, D. Rational approximations to algebraic numbers\index{rational approximations to algebraic numbers}\index{algebraic number}, \textit{Mathematika}, 4 (1957), 125-31.
    \item \hspace{1.5em} The p-adic generalization of the Thue\index{Thue, A.}-Siegel\index{Siegel, C. L.}-Roth\index{Roth, K. F.} theorem, \textit{ibid.} 5 (1958), 40-8.
    \item ROTH, K. F. Rational approximations to algebraic numbers\index{rational approximations to algebraic numbers}\index{algebraic number}, \textit{Mathematika}, 2 (1955), 1-20; corrigendum, \textit{ibid}. 168.
    \item Schinzel\index{Schinzel, A.}, A. On two theorems of Gelfond\index{Gelfond, A. O.} and some of their applications, \textit{Acta Arith.} 13 (1968), 177-236.
    \item \hspace{1.5em} An improvement on Runge's theorem on Diophantine equations\index{Diophantine equations}, Comment. Pontificia Acad. Sci. 2 (1969), no. 20, 1-9.
    \item \hspace{1.5em} Primitive divisors of the expression  $A^n-B^n$  in algebraic number\index{algebraic number} fields, J.M. 269 (1974), 27-33.
    \item SCHMIDT, W. M. Simultaneous approximation and algebraic independence of numbers\index{numbers}, \textit{Bull. Amer. Math. Soc.} 68 (1962), 475-8.
    \item \hspace{1.5em} Uber simultane Approximation algebraischer Zahlen durch rationale, Acta Math. 114 (1965), 159–206.
    \item \hspace{1.5em} Simultaneous approximation to a basis of a real number field, Amer. J. Math. 88 (1966), 517-27.
    \item \hspace{1.5em} Some diophantine equations\index{Diophantine equations} in three variables with only finitely many solutions, \textit{Mathematika}, 14 (1967), 113-20.
    \item \hspace{1.5em} T-numbers\index{T-number}\index{numbers} do exist, Symposia Math. IV, INDAM, Rome, 1968 (Academic Press, London, 1970), pp. 3-26.
    \item \hspace{1.5em} Simultaneous approximation to algebraic numbers\index{algebraic number} by rationals, \textit{Acta Math.} 125 (1970), 189-201.
    \item \hspace{1.5em} Mahler\index{Mahler, K.}'s T-numbers\index{T-number}, Proc. Symposia Pure Math., vol. 20 (Amer. Math. Soc. 1971), pp. 275-86.
    \item \hspace{1.5em} Linearformen mit algebraischen Koeffizienten I, II, J. Number Th. 3 (1971), 253-77; M.A. 191 (1971), 1-20.
    \item \hspace{1.5em} Norm form equations\index{norm form equation}, Ann. Math. 96 (1972), 526-51.
    \item Schneider\index{Schneider, Th.}, Th. Transzendenzuntersuchungen periodischer Funktionen: I Transzendenz von Potenzen; II Transzendenzeigenschaften elliptischer Funktionen, J.M. 172 (1934), 65-74.
    \item Über die Approximation algebraischer Zahlen, J.M. 175 (1936), 182-92.
    \item Arithmetische Untersuchungen elliptischer Integrale, M.A. 113 (1937), 1-13.
    \item Zur Theorie der Abelschen Funktionen und Integrale. I M. 183 (1941)
    \item Zur Theorie der Abelschen Funktionen und Integrale, J.M. 183 (1941), 110-28.
    \item Über eine Dysonsche Verschärfung des Siegel\index{Siegel, C. L.}-Thueschen Satzes, Arch. Math. 1 (1949), 288-95.
    \item Ein Satz über ganzwertige Funktionen als Prinzip für Transzendenzbeweise, M.A. 121 (1949), 131-40.
    \item Zur Annäherung der algebraischen Zahlen durch rationale, J.M. 188 (1950), 115–28.
    \item Anwendung eines abgeänderten Roth\index{Roth, K. F.}-Ridoutschen Satzes auf diophantische Gleichungen, M.A. 169 (1967), 177-82.
    \item Über p-adische Kettenbrüche, Symposia Math. IV, INDAM, Rome, 1968 (Academic Press, London, 1970), pp. 181-9.
    \item Schwarz\index{Schwarz, W.}, W. Remarks on the irrationality and transcendence of certain series, \textit{Math. Scand.} 20 (1967), 269-74.
    \item SHIDLOVSKY, A. B. On bounds for the measure of transcendence\index{measure of transcendence} of a subclass of the numbers\index{numbers}  $\alpha^{\beta}$ , V.M. 6 (1951), 17-28.
    \item \hspace{1.5em} On the transcendence and algebraic independence of the values of integral functions of some classes, and similar titles, \textit{D.A.N.} 96 (1954), 697-700; 100 (1955), 221-4; 103 (1955), 979-80; 105 (1955), 35-7; 108 (1956), 400-3.
    \item \hspace{1.5em} On a criterion for algebraic independence of the values of a class of integral functions, I.A.N. 23 (1959), 35-66; = A.M.S.T. 22 (1962), 339-70.
    \item \hspace{1.5em} On the transcendence and algebraic independence of some E-functions\index{E-functions}, and similar titles, V.M. 5 (1960), 19–28; 5 (1961), 44–59; I.A.N. 26 (1962), 877–910; = A.M.S.T. 50 (1966), 141–77.
    \item \hspace{1.5em} On a generalization of Lindemann\index{Lindemann, F.}'s theorem\index{Lindemann's theorem}, D.A.N. 138 (1961), 1301-4; = 2 (1961), 841-4.
    \item \hspace{1.5em} A general theorem on the algebraic independence of \textit{E}-functions, and similar titles, \textit{D.A.N.} 169 (1966), 42-5; 171 (1966), 810-13; \textit{Litovsk Mat. Sb.} 1 (1966), 129-30.
    \item \hspace{1.5em} On estimates of the measure of transcendence\index{measure of transcendence} of E-functions\index{E-functions}, U.M. 22 (1967), 245-56; Mat. Zametki, 2 (1967), 33-44.
    \item \hspace{1.5em} Algebraic independence of the values of certain hypergeometric E-functions\index{E-functions}, Trudy Moskov, 18 (1968), 55-64; = 18 (1968), 59-69.
    \item \hspace{1.5em} Transcendence and algebraic independence of values of E-functions\index{E-functions}, Proc. Internat. Congress Math. (Moscow, 1968), pp. 299-307.
    \item \hspace{1.5em} A certain theorem of C. Siegel\index{Siegel, C. L.}, V.M. 24 (1969), 39-42.
    \item SHOREY, T. N. On a theorem of Ramachandra\index{Ramachandra, K.}, Acta Arith. 20 (1972), 215-21. Linear forms in the logarithms of algebraic numbers\index{logarithms of algebraic numbers}\index{algebraic number} with small coefficients I, II, J. Indian Math. Soc. 38 (1974), 271-92.
    \item \hspace{1.5em} Algebraic independence of certain numbers\index{numbers} in the p-adic domain, I.M. 34 (1972), 423-35, and related work, ibid. 436-42.
    \item \hspace{1.5em} On gaps between numbers\index{numbers} with a large prime factor II, Acta Arith. 25 (1974), 365-73.
    \item SIEGEL, C. L. Approximation algebraischer Zahlen, M.Z. 10 (1921), 173–213. Über Näherungswerte algebraischer Zahlen, M.A. 84 (1921), 80–99.
    \item \hspace{1.5em} Über den Thueschen Satz, Norske Vid. Selsk. Skr., no. 16 (1921).
    \item \hspace{1.5em} Uber einige Anwendungen diophantischer Approximationen, Abh. Preuss. Akad. Wiss. no. 1 (1929).
    \item \hspace{1.5em} Über die Perioden elliptischer Funktionen, J.M. 167 (1932), 62-9.
    \item SIEGEL, C. L. The integer solutions of the equation  $y^2 = ax^n + bx^{n-1} + ... + k$ , J. London Math. Soc. 1 (1926), 66-8 (under the pseudonym X).
    \item \hspace{1.5em} Die Gleichung  $ax^n by^n = c$ , M.A. 114 (1937), 57-68.
    \item \hspace{1.5em} Abschatzung von Einheiten, Göttinger Nachr. (1969), 71-86.
    \item \hspace{1.5em} Einige Erläuterungen zu Thues Untersuchungen über Annäherungswerte algebraischer Zahlen und diophantische Gleichungen, ibid. (1970), 169-95.
    \item SLESORAITENE, R. The Mahler\index{Mahler, K.}-Sprindžuk\index{Sprindzhuk, V. G.} theorem for polynomials of the third degree in two variables, and similar titles, \textit{Litovsk. Mat. Sb.} 9 (1969), 627-34; 10 (1970), 367-74, 545-64, 791-813; 13 (1973), 177-88.
    \item ŠMELEV\index{Smelev, A. A.}, A. A. The algebraic independence of a certain class of transcendental numbers\index{transcendental number}\index{numbers}, and similar titles, \textit{Mat. Zametki}, 3 (1968), 51-8; 4 (1968), 341-8, 525-32; 5 (1969), 117-28; 7 (1970), 203-10.
    \item \hspace{1.5em} Algebraic independence of values of certain E-functions\index{E-functions}, Izv. Vysš Učebn. Zaved. Matematika (1969), no. 4 (83), 103-12.
    \item SPIRA, R. A lemma in transcendental number\index{transcendental number} theory, Trans. Amer. Math. Soc. 146 (1949), 457-64 (but see review no. 5307 in Math. Reviews, 41).
    \item Sprindžuk\index{Sprindzhuk, V. G.}\index{Sprindžuk, V. G.}, V. G. On some general problems of approximating numbers by algebraic numbers\index{algebraic number}, \textit{Litovsk Mat. Sb.} 2 (1962), 129-45.
    \item \hspace{1.5em} On a classification of transcendental numbers\index{transcendental number}\index{numbers}, ibid. 2 (1962), 215-9.
    \item \hspace{1.5em} A proof of Mahler\index{Mahler, K.}'s conjecture\index{Mahler's conjecture} on the measure of the set of S-numbers\index{S-number}\index{S*-number}\index{numbers}, and similar titles, ibid. 2 (1962), 221-6; D.A.N. 154 (1964), 783-6; 155 (1964), 54-6; = 5 (1964), 183-7, 361-3; U.M. 19 (1964), 191-4; I.A.N. 29 (1965), 379-436; = A.M.S.T. 51 (1966), 215-72.
    \item \hspace{1.5em} A metric theorem on the smallest integral polynomials of several variables, Dokl. Akad. Nauk BSSR, 11 (1967), 5-6.
    \item \hspace{1.5em} Concerning Baker\index{Baker, R. C.}\index{Baker, A.}'s theorem on linear forms in logarithms\index{linear forms in logarithms}, \textit{ibid}. 11 (1967), 767-9.
    \item \hspace{1.5em} The finiteness of the number of rational and algebraic points on certain transcendental curves, D.A.N. 177 (1967), 524-7.
    \item \hspace{1.5em} Estimates of linear forms with p-adic logarithms of algebraic numbers\index{logarithms of algebraic numbers}\index{algebraic number}, Vesci Akad. Navuk BSSR, Ser. Fiz-Mat. (1968), no. 4, 5-14.
    \item \hspace{1.5em} Irrationality of the values of certain transcendental functions, I.A.N. 32 (1968), 93-107.
    \item \hspace{1.5em} Effectivization in certain problems of Diophantine approximation theory, Dokl. Akad. Nauk BSSR, 12 (1968), 293-7.
    \item \hspace{1.5em} On the metric theory of 'nonlinear' diophantine approximations, \textit{ibid.} 13 (1969), 298-301.
    \item \hspace{1.5em} On the theory of the hypergeometric functions\index{hypergeometric function} of Siegel\index{Siegel, C. L.}, \textit{ibid.} 13 (1969), 389-91.
    \item \hspace{1.5em} Effective estimates in 'ternary' exponential Diophantine equations\index{Diophantine equations}, \textit{ibid}. 13 (1969), 777-80.
    \item \hspace{1.5em} Effective bounds on rational approximations to algebraic numbers\index{rational approximations to algebraic numbers}\index{algebraic number}, \textit{ibid.} \textbf{14} (1970), 681-4, and similar titles, \textit{ibid.} \textbf{15} (1971), 101-4; \textit{I.A.N.} \textbf{35} (1971), 991-1007.
    \item \hspace{1.5em} A new application of p-adic analysis to representations of numbers\index{numbers} by binary forms, I.A.N. 34 (1970), 1038-63.
    \item \hspace{1.5em} New applications of analytic and p-adic methods in Diophantine approximations, \textit{Proc. Internat. Congress Math.} Nice, 1970, vol. 1 (Gauthier-Villars, Paris, 1971).
    \item \hspace{1.5em} On the largest prime factor of a binary form, Dokl. Akad. Nauk BSSR, 15 (1971), 389-91, and related work 17 (1973), 685-8.
    \item \hspace{1.5em} On bounds for the units in algebraic number\index{algebraic number} fields, ibid. 15 (1971), 1065-8.
    \item On bounds for the solutions of the Thue\index{Thue, A.} equation, I.A.N. 36 (1972), 712-41. The method of trigonometrical sums in the metrical theory\index{metrical theory} of Diophantine approximation, Trudy Mat. Inst. Steklov, 128 (1972), 212-28.
    \item STARK, H. M. A complete determination of the complex quadratic fields with class-number one, \textit{Michigan Math. J.} 14 (1967), 1–27.
    \item \hspace{1.5em} On the 'gap' in a theorem of Heegner\index{Heegner, K.}, J. Number Th. 1 (1969), 16-27.
    \item \hspace{1.5em} A historical note on complex quadratic fields with class-number one, \textit{Proc. Amer. Math. Soc.} 21 (1969), 254-5.
    \item \hspace{1.5em} A transcendence theorem for class number problems I, II, Ann. Math. 94 (1971), 153-73; 96 (1972), 174-209.
    \item \hspace{1.5em} Effective estimates of solutions of some diophantine equations\index{Diophantine equations}, Acta Arith. 24 (1973), 251-9.
    \item \hspace{1.5em} Further advances in the theory of linear forms in logarithms\index{linear forms in logarithms}, \textit{Diophantine} approximation and its applications (Academic Press, London, 1973), pp. 255-93.
    \item STEPANOV, S. A. The approximation of an algebraic number\index{algebraic number} by algebraic numbers\index{algebraic number} of a special form, V.M. 22 (1967), 78-86.
    \item STRAUS, E.G. On entire functions with algebraic derivatives at certain algebraic points, Ann. Math. 52 (1950), 188-98.
    \item Tartakovsky\index{Tartakovsky, V. A.}, V. A. A uniform estimate for the number of representations of unity by a binary form of degree  $n \ge 3$ , D.A.N. 193 (1970), 764; = 11 (1970), 1026-7.
    \item Thue\index{Thue, A.}, A. Bemerkungen über gewisse Näherungsbrüche algebraischer Zahlen, Norske Vid. Selsk. Skr. No. 3 (1908).
    \item \hspace{1.5em} Über rationale Annäherungswerte der reellen Wurzel der ganzen Funktion dritten Grades  $x^3 ax b$ , \textit{ibid}. no. 6 (1908) (see also no. 7).
    \item \hspace{1.5em} Über Annäherungswerte algebraischer Zahlen, J.M. 135 (1909), 284-305.
    \item \hspace{1.5em} Ein Fundamentaltheorem zur Bestimmung von Annäherungswerten aller Wurzeln gewisser ganzer Funktionen, J.M. 138 (1910), 96-108.
    \item \hspace{1.5em} Über eine Eigenschaft, die kiene transzendente Grösse haben kann, Norske Vid. Selsk. Skr. no. 20 (1912).
    \item \hspace{1.5em} Berechnung aller Lösungen gewisser Gleichungen von der Form  $ax^r by^r = f$ , \textit{ibid}. no. 4 (1918).
    \item TIJDEMAN, R. On the number of zeros of general exponential polynomials\index{exponential polynomials}, I.M. 33 (1971), 1-7.
    \item \hspace{1.5em} On the algebraic independence of certain numbers\index{numbers}, I.M. 33 (1971), 146-62. On the maximal distance of numbers\index{numbers} with a large prime factor, J. London Math. Soc. 5 (1972), 313-20.
    \item \hspace{1.5em} An auxiliary result in the theory of transcendental numbers\index{transcendental number}\index{numbers}, J. Number Th. 5 (1973), 80-94.
    \item \hspace{1.5em} On integers with many small prime factors, Compositio Math. 26 (1973), 319-30.
    \item UCHIYAMA, S. On the Thue\index{Thue, A.}-Siegel\index{Siegel, C. L.}-Roth\index{Roth, K. F.} theorem I, II, III, Proc. Japan Acad. 35 (1959), 413-6, 525-9; 36 (1960), 1-2.
    \item VELDRAMP, G. R. Ein Transzendenz-Satz für p-adische Zahlen, J. London Math. Soc. 15 (1940), 183-92.
    \item VINOGRADOV, A. I. and SPRINDŽUK\index{Sprindzhuk, V. G.}, V. G. The representation of numbers\index{numbers} by binary forms, \textit{Mat. Zametki}, 3 (1968), 369-76.
    \item VOLEMANN, B. Zur kubischen Fall der Mahlerschen Vermutung, M.A. 139 (1959), 87-90.
    \item \hspace{1.5em} Zur Mahlerschen Vermutung im Komplexen, M.A. 140 (1960), 351-9.
    \item \hspace{1.5em} Ein metrischer Beitrag über Mahlersche S-Zahlen I, J.M. 203 (1960), 154-6.
    \item VOLEMANN, B. The real cubic case of Mahler\index{Mahler, K.}'s conjecture\index{Mahler's conjecture}, \textit{Mathematika}, 8 (1961), 55-7.
    \item \hspace{1.5em} Zur metrischen Theorie der S-Zahlen I, II, J.M. 209 (1962), 201-10; 213 (1964), 58-65.
    \item WALDSCHMIDT, M. Solution d'un problème de Schneider\index{Schneider, Th.} sur les nombres transcendants, C.R. 271 (1970), 697-700.
    \item \hspace{1.5em} Amélioration d'un théorème de Lang\index{Lang, S.} sur l'indépendance algébrique d'exponentielles, C.R. 272 (1971), 413-5.
    \item \hspace{1.5em} Indépendance algébrique des valeurs de la fonction exponentielle, Bull. Soc. Math. France, 99 (1971), 285-304.
    \item \hspace{1.5em} Propriétés arithmétiques des valeurs de fonctions méromorphes algébriquement indépendantes, \textit{Acta Arith.} 23 (1973), 19-88.
    \item \hspace{1.5em} Solution du huitième probleme de Schneider\index{Schneider, Th.}, J. Number Th. 5 (1973), 191-202.
    \item WALLISSER, R. Zur Approximation algebraischer Zahlen durch arithmetisch charakterisierte algebraische Zahlen, Arch. Math. (Basel), 20 (1969), 384– 91.
    \item \hspace{1.5em} Zur Transzendenz der Werte der Exponentialfunktion, Monatsh. Math. 73 (1969), 449-60.
    \item \hspace{1.5em} Über Produkte transzendenter Zahlen, J.M. 258 (1973), 62-78.
    \item WEIERSTRASS, K. Zu Lindemann\index{Lindemann, F.}'s Abhandlung 'Uber die Ludolph'sche Zahl', S.-B. Preuss Akad. Wiss. (1885), 1067-85; = Werke II, 341-62.
    \item Wirsing\index{Wirsing, E.}, E. Approximation mit algebraischen Zahlen beschränkten Grades, J.M. 206 (1961), 67-77.
    \item \hspace{1.5em} On approximations of algebraic numbers\index{algebraic number} by algebraic numbers\index{algebraic number} of bounded degree, \textit{Proc. Symposia Pure Math.} vol. 20 (Amer. Math. Soc. 1971), pp. 213-47.
\end{hanglist}
